% 08-expressions.tex — Clause 8: Expressions
\chapter{Expressions}
\label{sec:08-expressions}
\index{expression}

\pnum
An expression is a syntactic construct that evaluates to a value.  Every
expression has a type, determined statically at compile time.  The grammar
production for \gramref{expression} is in Annex~A.

\pnum
Expressions are evaluated strictly left-to-right: for any binary operator
$a\ \mathit{op}\ b$ or function call \lain{f(a, b)}, the left operand
(or earlier argument) is fully evaluated before the right operand (or
later argument).

\section{Operator precedence and associativity}
\label{sec:expr.precedence}
\index{operator precedence}

\pnum
The following table lists all binary operators in order of decreasing
precedence.  All operators listed on the same row have equal precedence.

\begin{longtable}{@{}llll@{}}
\toprule
\textbf{Prec.} & \textbf{Operators} & \textbf{Assoc.} & \textbf{Result type} \\
\midrule
\endfirsthead
\toprule
\textbf{Prec.} & \textbf{Operators} & \textbf{Assoc.} & \textbf{Result type} \\
\midrule
\endhead
\bottomrule
\endlastfoot
1 (high) & \lain{()} \lain{[]} \lain{.}              & left  & (context-dependent) \\
2        & unary \lain{-} \lain{!} \lain{not} \lain{mov} \lain{ref} \lain{mut} & right & (operand type) \\
3        & \lain{as}                                  & left  & target type \\
4        & \lain{*} \lain{/} \lain{\%}                & left  & wider operand type \\
5        & \lain{+} \lain{-}                          & left  & wider operand type \\
6        & \lain{<<} \lain{>>}                        & left  & left operand type \\
7        & \lain{\&}                                  & left  & wider operand type \\
8        & \texttt{\textasciicircum}                                & left  & wider operand type \\
9        & \lain{|}                                   & left  & wider operand type \\
10       & \lain{<} \lain{>} \lain{<=} \lain{>=} \lain{in} & left & \lain{bool} \\
11       & \lain{==} \lain{!=}                        & left  & \lain{bool} \\
12       & \lain{and}                                 & left  & \lain{bool} \\
13 (low) & \lain{or}                                  & left  & \lain{bool} \\
\end{longtable}

\pnum
Parentheses override precedence.

\section{Primary expressions}
\label{sec:expr.primary}
\index{primary expression}

\pnum
A primary expression is one of: an identifier, a literal, a parenthesized
expression, a block expression, an \lain{if} expression, a \lain{case}
expression, a constructor expression, or a \lain{comptime} expression.

\pnum
\textbf{Identifier expression}: evaluates to the current value of the named
variable.

\begin{constraint}
  The identifier shall be declared in scope
  (\diagcode{E013}).
\end{constraint}

\begin{constraint}
  The identifier shall have been initialized before use
  (\diagcode{E005}).
\end{constraint}

\section{Arithmetic operators}
\label{sec:expr.arith}
\index{arithmetic operator}

\pnum
The binary arithmetic operators are \lain{+} (addition), \lain{-}
(subtraction), \lain{*} (multiplication), \lain{/} (division), and
\lain{\%} (remainder).  The result type is the wider of the two operand
types (\S\,\ref{sec:types.rank}).

\pnum
Integer division truncates toward zero.  \lain{f32}/\lain{f64} division
follows IEEE~754.

\begin{constraint}
  Integer division or remainder where the divisor is statically provable to
  be zero (via value range analysis, \S\,\ref{sec:13-vra}) is
  \illformed{}~(\diagcode{E015}).
\end{constraint}

\begin{constraint}
  \lain{\%} shall not be applied to floating-point operands.
\end{constraint}

\pnum
The unary \lain{-} operator negates its operand.  For unsigned integer
types the result wraps modulo $2^N$.

\section{Comparison operators}
\label{sec:expr.compare}
\index{comparison operator}

\pnum
The comparison operators \lain{==}, \lain{!=}, \lain{<}, \lain{>},
\lain{<=}, \lain{>=} compare two values and yield \lain{bool}.

\begin{constraint}
  Operands of a comparison operator shall have compatible types
  (\S\,\ref{sec:types.compat}).
\end{constraint}

\begin{constraint}
  \lain{==} and \lain{!=} shall not be applied to struct types, array types,
  or enum types (\S\,\ref{sec:types.enum}).
\end{constraint}

\section{Logical operators}
\label{sec:expr.logical}
\index{logical operator}

\pnum
\lain{and} and \lain{or} are short-circuit operators: the right operand
is not evaluated if the result is determined by the left operand alone.
\lain{not} and \lain{!} are prefix unary negation.  All yield \lain{bool}.

\begin{constraint}
  Operands of \lain{and}, \lain{or}, \lain{not}, and \lain{!} shall have
  type \lain{bool}.
\end{constraint}

\section{Bitwise operators}
\label{sec:expr.bitwise}
\index{bitwise operator}

\pnum
The bitwise operators \lain{\&}, \lain{|}, \texttt{\textasciicircum}, \lain{<<}, \lain{>>}
operate on integer operands.  The result type is the wider operand type for
\lain{\&}, \lain{|}, \texttt{\textasciicircum}; for shift operators, the result type is
the left operand type.

\begin{constraint}
  Operands of bitwise operators shall be integer types.
\end{constraint}

\begin{constraint}
  The shift amount (right operand of \lain{<<} and \lain{>>}) shall
  satisfy $0 \leq \mathit{shift} < N$ where $N$ is the bit width of the
  left operand.  A shift amount outside this range is
  \illformed{}~(\diagcode{E100}).
\end{constraint}

\section{The \texttt{in} operator}
\label{sec:expr.in}
\index{in operator}
\index{in guard}

\pnum
The expression \lain{idx in arr} evaluates to \lain{true} if and only if
$0 \leq \mathit{idx} < \mathit{arr.len}$.  The operand \textit{arr} shall
be an array or slice type; \textit{idx} shall be an integer type.

\pnum
When \lain{idx in arr} appears as the condition of a \lain{while} or
\lain{if} statement (or as the leftmost operand of an \lain{and}-chain
condition), it establishes an \textbf{in guard}: inside the body of that
statement, the implementation shall accept \lain{arr[idx]} without
requiring further bounds verification.

\pnum
In-guards propagate through \lain{and}-chains: if the left operand of
\lain{and} establishes an in-guard, that guard is active when evaluating
the right operand.

\begin{example}
\begin{laincode}
// In-guard: arr[i] is safe without explicit VRA
while i in arr {
    var x = arr[i]   // OK: in-guarded
    i = i + 1
}

// And-chain propagation:
while l.pos in l.src and (l.src[l.pos] as int) != '"' {
    l.pos = l.pos + 1
}
\end{laincode}
\end{example}

\begin{constraint}
  An in-guard of the form \lain{idx in arr} covers only \lain{arr[idx]}
  exactly.  Accesses \lain{arr[idx + k]} for any $k \neq 0$ are not
  covered by this guard and require separate verification or an
  \lain{unsafe} block.
\end{constraint}

\section{Cast expression (\texttt{as})}
\label{sec:expr.cast}
\index{cast expression}
\index{as operator}

\pnum
The expression \lain{expr as T} explicitly converts the value of
\textit{expr} to type \lain{T}.  The following casts are permitted:

\begin{longtable}{@{}lll@{}}
\toprule
\textbf{From} & \textbf{To} & \textbf{Semantics} \\
\midrule
\endfirsthead
\toprule
\textbf{From} & \textbf{To} & \textbf{Semantics} \\
\midrule
\endhead
\bottomrule
\endlastfoot
integer & integer & Truncation (narrowing) or extension (widening) \\
integer & float   & Nearest representable value (IEEE~754) \\
float   & integer & Truncation toward zero \\
float   & float   & Precision change \\
pointer & pointer & Reinterpretation (\lain{unsafe} required) \\
integer & pointer & Reinterpretation (\lain{unsafe} required) \\
pointer & integer & Reinterpretation (\lain{unsafe} required) \\
\end{longtable}

\begin{constraint}
  Casts between pointer types, between a pointer and an integer, or between
  any type and a pointer shall occur only inside an \lain{unsafe} block.
  Such a cast outside \lain{unsafe} is \illformed{}~(\diagcode{E012}).
\end{constraint}

\begin{constraint}
  A cast between non-numeric, non-pointer types (e.g., struct to integer)
  is \illformed{}.
\end{constraint}

\section{Move expression (\texttt{mov})}
\label{sec:expr.mov}
\index{move expression}

\pnum
\lain{mov expr} transfers ownership of the value produced by \textit{expr}.
The source is placed in the Consumed state; any subsequent use of the
source variable is \illformed{}~(\diagcode{E001}).  See \S\,\ref{sec:11-ownership}.

\section{Borrow expressions (\texttt{ref}, \texttt{mut})}
\label{sec:expr.borrow}
\index{borrow expression}

\pnum
\lain{ref expr} creates a shared (read-only) borrow of \textit{expr}.
\lain{mut expr} creates a mutable borrow of \textit{expr}.
Both expressions are evaluated in argument contexts to pass values without
transferring ownership.  See \S\,\ref{sec:11-ownership}.

\section{Field access}
\label{sec:expr.field}
\index{field access}

\pnum
The expression \lain{expr.field} accesses the named field of a struct or
ADT value, or a property of an array or slice (\lain{.len}, \lain{.data}).

\begin{constraint}
  The field name shall exist on the type of \textit{expr}.
\end{constraint}

\begin{constraint}
  Writing to a field (\lain{expr.field = v}) requires \textit{expr} to be
  a mutable lvalue.
\end{constraint}

\section{Index expression}
\label{sec:expr.index}
\index{index expression}

\pnum
The expression \lain{arr[idx]} accesses the element of array or slice
\textit{arr} at integer index \textit{idx}.

\begin{constraint}
  The value range analysis shall prove $0 \leq \mathit{idx} < \mathit{arr.len}$
  at every index expression.  An index expression for which this cannot be
  proved is \illformed{}~(\diagcode{E014}).  An \lain{in} guard
  (\S\,\ref{sec:expr.in}) or an \lain{unsafe} block suppresses this
  requirement.
\end{constraint}

\section{Function call}
\label{sec:expr.call}
\index{function call}

\pnum
A function call expression \lain{f(args)} evaluates the arguments and
transfers control to the callee.  Arguments are evaluated left-to-right.

\pnum
Argument ownership modes shall be given explicitly at the call site:

\begin{longtable}{@{}ll@{}}
\toprule
\textbf{Syntax} & \textbf{Meaning} \\
\midrule
\endfirsthead
\toprule
\textbf{Syntax} & \textbf{Meaning} \\
\midrule
\endhead
\bottomrule
\endlastfoot
\lain{f(x)}     & shared borrow (default; parameter mode must be \lain{shared}) \\
\lain{f(var x)} & mutable borrow (parameter mode must be \lain{var}) \\
\lain{f(mov x)} & ownership transfer (parameter mode must be owned/\lain{mov}) \\
\end{longtable}

\begin{constraint}
  The mode annotation at the call site shall match the declared parameter
  mode.  A mismatch is \illformed{}.
\end{constraint}

\begin{constraint}
  The number of arguments shall equal the number of parameters (unless the
  callee is a variadic C function declared with \lain{...}).
\end{constraint}

\pnum
\textbf{Two-phase borrows}: during argument evaluation, borrows of the
receiver of a method call are placed in the RESERVED phase and promoted to
ACTIVE after all arguments are evaluated
(see \S\,\ref{sec:11-ownership}).

\section{Block expression}
\label{sec:expr.block}
\index{block expression}

\pnum
A block expression \texttt{\{ stmts... expr \}} evaluates each statement
in order, then evaluates the final expression, whose value is the value of
the block.  If the final expression is absent, the block has no value.

\section{If expression}
\label{sec:expr.if}
\index{if expression}

\pnum
An \lain{if} expression of the form
\lain{if cond \{ e1 \} else \{ e2 \}} evaluates \textit{cond}; if true,
evaluates $e_1$; otherwise, evaluates $e_2$.

\begin{constraint}
  The types of $e_1$ and $e_2$ shall be compatible.  If they are not, the
  expression is \illformed{}.
\end{constraint}

\begin{constraint}
  \textit{cond} shall have type \lain{bool}.
\end{constraint}

\section{Match expression (\texttt{case})}
\label{sec:expr.match}
\index{case expression}

\pnum
A \lain{case} expression (\gramref{match-expression}) selects one of
several alternatives based on pattern matching.  All arms shall produce
values of compatible types; the type of the \lain{case} expression is that
common type.  See \S\,\ref{sec:15-match} for full exhaustiveness rules.

\section{Constructor expression}
\label{sec:expr.constructor}
\index{constructor expression}

\pnum
A constructor expression builds a value of a struct or ADT type by
supplying all field values.  See \S\,\ref{sec:types.struct} and
\S\,\ref{sec:types.adt} for constructor rules.

\section{Compound assignment}
\label{sec:expr.compound-assign}
\index{compound assignment}

\pnum
Compound assignment \lain{x op= e} is syntactic sugar for
\lain{x = x op e}.  The left-hand side shall be a mutable lvalue.
All compound forms of the arithmetic and bitwise operators are supported
(see \S\,\ref{sec:lex.operators}).
